\question[3]
Erkläre das Prinzip der Namensauflösung durch DNS, indem du es mit einer 
Kontaktliste auf einem Smartphone vergleichst. 
Gehe darauf ein, was in beiden Fällen der \enquote{Name} und was die \enquote{Adresse} ist.

\begin{solutionbox}{4cm}

In einer Kontaktliste tippt man einen Namen ein
und das Smartphone wählt automatisch die passende Rufnummer,
ohne dass man sie selbst kennen oder eintippen muss.

Beim DNS (Domain Name System) gibt man einen gut merkbaren
Domainnamen (z. B. \texttt{www.example.de}) ein, und der DNS-Server
löst diesen Namen automatisch in die zugehörige, schwer zu
merkende IP-Adresse (z. B. \texttt{93.184.216.34}) auf.

\end{solutionbox}